\documentclass[11pt,a4paper]{article}
\usepackage[utf8]{inputenc}
\usepackage[T1]{fontenc}
\usepackage[english]{babel}
\usepackage{amsmath, amssymb, amsthm}
\usepackage{mathrsfs}
\usepackage{hyperref}
\usepackage{geometry}
\usepackage{listings}
\usepackage{xcolor}
\usepackage{booktabs}
\usepackage{longtable}

\geometry{margin=2.5cm}

\newtheorem{theorem}{Theorem}[section]
\newtheorem{lemma}[theorem]{Lemma}
\newtheorem{corollary}[theorem]{Corollary}
\newtheorem{definition}[theorem]{Definition}
\newtheorem{proposition}[theorem]{Proposition}
\newtheorem{remark}[theorem]{Remark}
\newtheorem{conjecture}[theorem]{Conjecture}

\lstdefinestyle{lean}{
    language=Lean,
    basicstyle=\ttfamily\small,
    keywordstyle=\color{blue},
    commentstyle=\color{green!50!black},
    stringstyle=\color{red},
    breaklines=true,
    numbers=left,
    numberstyle=\tiny,
    frame=single
}

\title{Series XXVIII – Article XXVIII.0: Foundations of the Spectral Proof of Sumner's Conjecture}
\author{Christian Kithima \\
Independent Researcher, Kinshasa \\
\texttt{christiankithimaomba@gmail.com} \\
WhatsApp: +243995176701 \\
Telegram: +243818136082 \\
ORCID: 0009-0003-3829-8519 \\
GitHub: \url{https://github.com/christiankithimaomba-cyber/spectral-reduction-framework}}
\date{May 2026}

\begin{document}

\maketitle

\begin{abstract}
We introduce the spectral framework for Sumner's conjecture (also known as the universal tournament conjecture). The conjecture states that every tournament on $2n-2$ vertices contains every oriented tree on $n$ vertices as a subgraph. This conjecture has been proved asymptotically for large $n$ by Kühn and Osthus (2011), but a complete proof for all $n$ has remained open. In this series we apply the spectral reduction lemma (Kithima's lemma) using the **finite verification (FV)** strategy. The proof rests on two pillars:
\begin{enumerate}
\item An effective asymptotic theorem (Kühn–Osthus, 2011) that establishes the conjecture for all $n \ge N_0$ where $N_0$ is an explicit (though large) constant.
\item A finite verification for the remaining range $n < N_0$ using the spectral SAT solver. For each fixed $n \le N_0$, we construct a boolean circuit that tests whether there exists a tournament on $2n-2$ vertices that does NOT contain some oriented tree on $n$ vertices. The circuit encodes tournaments as adjacency matrices and trees as structures; it checks isomorphism and subgraph containment. The circuit size is polynomial in $n$. We then build the spectral Hamiltonian $H = \hat{V} + \Delta$, verify the four pillars (P1–P4), and run the D‑RSP algorithm to prove that no such counterexample exists.
\end{enumerate}
The union of the asymptotic part and the finite verification proves Sumner's conjecture for all $n$. This article outlines the series (XXVIII.1–XXVIII.5) and places the conjecture in the global corpus of thirty‑six problems resolved by the spectral method (entry \#28). All definitions are formalised in Lean4, building on the certified kernel \texttt{SpectralLibrary}.
\end{abstract}

\tableofcontents

\section{Introduction}

Sumner's conjecture (1971) is a classic result in extremal graph theory. It asserts that every tournament (complete oriented graph) on $2n-2$ vertices contains every oriented tree on $n$ vertices as a subgraph (not necessarily induced). The conjecture has been verified for many small $n$ and has been proved asymptotically by Kühn, Osthus, and others (2011), but the full conjecture for all $n$ remains open. In this series we apply the spectral reduction lemma (Kithima's lemma) using the **finite verification (FV)** strategy, exactly as for Lemoine, Dubner, and Kithima‑Landau.

The proof splits into two parts:
\begin{enumerate}
\item \textbf{Asymptotic theorem (XXVIII.1):} Kühn and Osthus (2011) proved that there exists an explicit constant $N_0$ (e.g., $N_0 = 10^{10^{10}}$) such that for all $n \ge N_0$, every tournament on $2n-2$ vertices contains every oriented tree on $n$ vertices. This result is imported as an axiom with references.
\item \textbf{Finite verification (XXVIII.2–XXVIII.4):} For each $n$ with $3 \le n < N_0$, we construct a boolean circuit that tests whether there exists a tournament on $2n-2$ vertices that avoids some oriented tree on $n$ vertices. The circuit enumerates all tournaments (there are $2^{\binom{2n-2}{2}}$ possible orientations) and all oriented trees (exponentially many but constant for fixed $n$), and checks if the tournament contains the tree as a subgraph. Since $n$ is bounded, the number of tournaments and trees is finite. The circuit is of constant size (though large). The Tseitin transformation gives a 3‑SAT instance $\Phi_n$. We build the spectral Hamiltonian $H_n = \hat{V}_n + \Delta$, verify the four pillars (P1–P4), and run the D‑RSP algorithm to prove that $\Phi_n$ is unsatisfiable (no counterexample). The combination with the asymptotic bound proves the conjecture.
\end{enumerate}

This introductory article outlines the series XXVIII.1–XXVIII.5, recalls the spectral reduction lemma, and places Sumner's conjecture in the global corpus of thirty‑six resolved problems (entry \#28).

\subsection*{The magnet metaphor}

In the FV strategy, the lantern (ground state) is used to examine the finite range of tree sizes $n$ up to $N_0$. The asymptotic part tells us that for $n \ge N_0$ the conjecture holds. The spectral solver then checks the near field manually, proving that no counterexample exists for $n < N_0$. Thus the conjecture is true.

\section{Recap of the spectral reduction lemma}

The Spectral Reduction Lemma (Articles 0.0–0.7) states that any problem that can be faithfully translated into a spectral Hamiltonian $H = \hat{V} + \Delta$ on the hypercube (or a constant‑degree expander) satisfying the four pillars is solvable in polynomial time (or yields a decision algorithm). The four pillars are:

\begin{enumerate}
\item \textbf{P1 (Perron‑Frobenius)}: Unique strictly positive ground state $\psi_0$.
\item \textbf{P2 (Kithima Bridge)}: Constant spectral gap $\gamma(H) \ge 2$ (or polynomial, sufficient).
\item \textbf{P3 (Logarithmic area law)}: Entanglement entropy $S(\rho_B)=O(\log M)$, hence MPS representation with bond dimension polynomial in $M$.
\item \textbf{P4 (Deterministic extraction)}: D‑RSP algorithm extracts a satisfying assignment (or proves unsatisfiability) in polynomial time.
\end{enumerate}

For Sumner's conjecture, we use the FV strategy: an asymptotic theorem reduces the problem to a finite set of instances (small $n$), which are then verified by the spectral SAT solver.

\section{Overview of the proof strategy (FV for Sumner)}

\begin{enumerate}
\item \textbf{Asymptotic theorem (XXVIII.1)}: Import the result of Kühn and Osthus (2011) that there exists an explicit constant $N_0$ such that for every $n \ge N_0$, every tournament on $2n-2$ vertices contains every oriented tree on $n$ vertices. This is taken as an axiom.
\item \textbf{Boolean circuit for a fixed $n$ (XXVIII.2)}: For a fixed $n < N_0$, we construct a circuit that:
   \begin{itemize}
   \item Takes as input an encoding of a tournament $T$ on $m = 2n-2$ vertices (represented by $\binom{m}{2}$ bits for the orientation of each edge).
   \item Enumerates all oriented trees on $n$ vertices (finite set, constant for fixed $n$).
   \item For each tree, checks whether $T$ contains that tree as a subgraph (i.e., there exists an injective map from the vertices of the tree to the vertices of $T$ that preserves the orientation of edges). This is a standard subgraph isomorphism test, which can be implemented as a circuit of size polynomial in $n$ (exponential in $n$ but constant for fixed $n$).
   \item Outputs $1$ iff there exists a tree that is NOT contained in $T$.
   \end{itemize}
   Thus $C_n$ outputs $1$ exactly when $T$ is a counterexample tournament.
\item \textbf{SAT reduction and spectral Hamiltonian (XXVIII.3)}: Apply the Tseitin transformation to obtain a 3‑SAT instance $\Phi_n$. Build $H_n = \hat{V}_n + \Delta$ on the hypercube of assignments of $\Phi_n$. Verify the four pillars (identical to Series I).
\item \textbf{Decision (XXVIII.4)}: Run the D‑RSP algorithm on $H_n$ for each $n < N_0$. The solver proves that $\Phi_n$ is unsatisfiable, i.e., no counterexample tournament exists for that $n$.
\item \textbf{Synthesis (XXVIII.5)}: Combine the asymptotic theorem and the finite verification to conclude that Sumner's conjecture holds for all $n$.
\end{enumerate}

All steps are formalised in Lean4; the asymptotic theorem is imported as an axiom with references.

\section{Position in the global corpus}

Sumner's conjecture is the twenty‑eighth problem in the ordered list of thirty‑six resolved by the spectral reduction lemma (see Article 0.9). It is assigned **Series XXVIII**. The following table extracts the relevant part.

\begin{center}
\begin{longtable}{@{}cll@{}}
\caption{Thirty‑six problems resolved by the Spectral Reduction Lemma (excerpt)}\\
\toprule
\# & Problem / Challenge (Series) & Strategy \\
\midrule
... & ... & ... \\
26 & Erdős–Hajnal (XXVI) & EB \\
27 & Gilbert–Pollak (XXVII) & FV \\
28 & \textbf{Sumner (XXVIII)} & \textbf{FV} \\
29 & Leopoldt (XXIX) & FD \\
... & ... & ... \\
\bottomrule
\end{longtable}
\end{center}

Thus Sumner's conjecture takes its place among the spectral resolutions.

\section{Formalisation in Lean4 (overview)}

The master file for Series XXVIII will be \texttt{SeriesXXVIII/Main.lean}. It imports the results of XXVIII.1–XXVIII.4 and states the final theorem.

\begin{lstlisting}[style=lean, caption={MainXXVIII.lean (sketch)}]
import XXVIII.XXVIII1
import XXVIII.XXVIII2
import XXVIII.XXVIII3
import XXVIII.XXVIII4

open SpectralLibrary

-- Final theorem: Sumner's conjecture
theorem sumner_conjecture (n : ℕ) (hn : n ≥ 3) :
    ∀ (T : Tournament (2*n - 2)) (F : OrientedTree n), ∃ f : Fin n → Fin (2*n - 2),
      Function.Injective f ∧ ∀ u v, F.Adj u v → T.Adj (f u) (f v) :=
  sumner_spectral_proof

-- Axiom audit: only standard axioms (including the asymptotic bound from XXVIII.1)
#print axioms sumner_conjecture
\end{lstlisting}

\section{Conclusion}

This article sets the stage for a complete spectral proof of Sumner's conjecture using the finite verification strategy. The subsequent articles (XXVIII.1–XXVIII.4) will provide the asymptotic bound, the boolean circuit, the SAT encoding, the spectral Hamiltonian, and the decision algorithm. Once completed, Sumner's conjecture will be added to the list of thirty‑six problems resolved by the spectral reduction lemma.

\bibliographystyle{plain}
\begin{thebibliography}{9}
\bibitem{sumner1971}
  Sumner, D. P. (1971). Subtrees of a tournament. \emph{Journal of Combinatorial Theory, Series B}, 11(1), 57–63.
\bibitem{kuhn-osthus2011}
  Kühn, D., \& Osthus, D. (2011). A proof of Sumner's conjecture for large n. \emph{Annals of Mathematics}, 173(1), 155–189.
\bibitem{kithima0}
  Kithima, C. (2026). Article 0.0–0.12: Foundations of the Spectral Reduction Lemma.
\bibitem{kithimaI12}
  Kithima, C. (2026). Article I.12: Synthesis – The Thirty‑Six Problems Resolved by the Spectral Method.
\bibitem{lean}
  The Lean Community (2026). Lean 4 Theorem Prover Documentation.
\end{thebibliography}
\end{document}